\section{Thiết kế node cảm biến}
Node cảm biến được xây dựng xung quanh Arduino Pro Mini. Các cảm biến TDS và pH được sử dụng để theo dõi chất lượng nước hoặc dung dịch chăm sóc cây. Cảm biến nhiệt độ, độ ẩm và CO2 phục vụ đánh giá chất lượng không khí trong căn hộ. Cảm biến ánh sáng giúp xác định điều kiện chiếu sáng cho cây và không gian sinh hoạt.

\subsection{Khối đọc cảm biến}
Chương trình trên node thực hiện chu kỳ đọc dữ liệu theo các bước: khởi tạo cảm biến, lấy mẫu nhiều lần, lọc nhiễu đơn giản bằng trung bình cộng, quy đổi ra đơn vị đo và lưu vào cấu trúc dữ liệu. Với cảm biến analog, giá trị ADC được chuyển đổi theo hệ số hiệu chuẩn. Với cảm biến digital, node đọc dữ liệu theo giao thức của module và kiểm tra lỗi trước khi sử dụng.

\subsection{Khối điều khiển}
Node điều khiển hai relay để bật/tắt quạt và đèn. Relay được điều khiển qua chân GPIO, có trạng thái mặc định an toàn khi node khởi động. Bộ phát IR được dùng để gửi mã tăng hoặc giảm nhiệt độ điều hòa. Khi nhận lệnh IR từ gateway, node phát chuỗi mã tương ứng và gửi phản hồi trạng thái để gateway cập nhật lên ThingsBoard.

\subsection{Định dạng gói tin node gửi lên gateway}
Gói tin uplink được thiết kế ngắn gọn để phù hợp với LoRa:
\begin{center}
\texttt{NODE\_ID,T,H,CO2,LUX,TDS,PH,FAN,LIGHT}
\end{center}
Trong đó \texttt{T} là nhiệt độ, \texttt{H} là độ ẩm, \texttt{CO2} là nồng độ CO2, \texttt{LUX} là ánh sáng, \texttt{TDS} là tổng chất rắn hòa tan, \texttt{PH} là chỉ số pH, \texttt{FAN} và \texttt{LIGHT} là trạng thái relay.

\subsection{Thuật toán đọc cảm biến}
Để giảm nhiễu, node không sử dụng trực tiếp một mẫu ADC đơn lẻ. Với các kênh analog, node lấy nhiều mẫu liên tiếp, loại bỏ một số mẫu bất thường nếu cần và tính trung bình. Với cảm biến digital, node kiểm tra mã lỗi hoặc điều kiện timeout trước khi cập nhật giá trị. Nếu một cảm biến lỗi trong một chu kỳ, node có thể giữ giá trị gần nhất và đánh dấu trạng thái lỗi để gateway không hiểu nhầm là dữ liệu mới hợp lệ.

\begin{lstlisting}[language=C,caption={Mã giả đọc và lọc cảm biến analog}]
float readAnalogAverage(uint8_t pin) {
    long sum = 0;
    for (int i = 0; i < 20; i++) {
        sum += analogRead(pin);
        delay(5);
    }
    float adc = sum / 20.0;
    return adc * 5.0 / 1023.0;
}
\end{lstlisting}

Sau khi có điện áp đo, firmware quy đổi thành giá trị vật lý. Các hệ số quy đổi cần được lưu trong cấu hình để thuận tiện hiệu chuẩn. Đối với pH, có thể dùng hai điểm chuẩn pH 4 và pH 7 hoặc pH 7 và pH 10. Đối với TDS, cần dùng dung dịch chuẩn và bù nhiệt độ nếu yêu cầu độ chính xác cao.

\subsection{Quản lý trạng thái node}
Node không chỉ gửi dữ liệu cảm biến mà còn gửi trạng thái thiết bị. Các trạng thái gồm RL1, RL2 và điều hòa. Việc gửi kèm trạng thái giúp app hiển thị đúng sau khi người dùng điều khiển hoặc sau khi node khởi động lại. Nếu app chỉ lưu trạng thái cục bộ, khi mất kết nối hoặc reset node, giao diện có thể không phản ánh đúng trạng thái thực tế.

\begin{lstlisting}[language=C,caption={Cấu trúc dữ liệu trạng thái node}]
struct NodeState {
    float temperature;
    float humidity;
    float co2;
    float lux;
    float tds;
    float ph;
    bool relay1;
    bool relay2;
    bool acPower;
    int acSetpoint;
};
\end{lstlisting}

\subsection{Xử lý lệnh điều khiển tại node}
Node xử lý lệnh theo nguyên tắc kiểm tra trước khi thực thi. Bản tin phải đúng loại \texttt{CMD}, đúng mã node và đúng tên thiết bị. Nếu lệnh không hợp lệ, node bỏ qua và không thay đổi trạng thái. Với relay, node đổi mức GPIO. Với điều hòa, node gọi hàm phát mã IR tương ứng.

\begin{lstlisting}[language=C,caption={Mã giả xử lý lệnh actuator}]
void handleCommand(Command cmd) {
    if (cmd.nodeId != NODE_ID) return;

    if (cmd.device == "RL1") {
        setRelay(PIN_RELAY_1, cmd.action == "ON");
    } else if (cmd.device == "RL2") {
        setRelay(PIN_RELAY_2, cmd.action == "ON");
    } else if (cmd.device == "AC") {
        sendIrCommand(cmd.action);
    }

    sendAck(cmd.device, cmd.action);
}
\end{lstlisting}
